\chapter{Advanced Compiler Use and Efficiency Hints}
\label{advanced-compiler}

\credits{by Robert MacLachlan}


\section{Advanced Compiler Introduction}

In \cmucl{}, as with any language on any computer, the path to efficient
code starts with good algorithms and sensible programming techniques,
but to avoid inefficiency pitfalls, you need to know some of this
implementation's quirks and features.  This chapter is mostly a fairly
long and detailed overview of what optimizations \python{} does.
Although there are the usual negative suggestions of inefficient
features to avoid, the main emphasis is on describing the things that
programmers can count on being efficient.

The optimizations described here can have the effect of speeding up
existing programs written in conventional styles, but the potential
for new programming styles that are clearer and less error-prone is at
least as significant.  For this reason, several sections end with a
discussion of the implications of these optimizations for programming
style.



\subsection{Types}

\python{}'s support for types is unusual in three major ways:
\begin{itemize}
  
\item Precise type checking encourages the specific use of type
  declarations as a form of run-time consistency checking.  This
  speeds development by localizing type errors and giving more
  meaningful error messages.  \xlref{precise-type-checks}.  \python{}
  produces completely safe code; optimized type checking maintains
  reasonable efficiency on conventional hardware
  (\pxlref{type-check-optimization}.)
  
\item Comprehensive support for the \clisp{} type system makes complex
  type specifiers useful.  Using type specifiers such as \code{or} and
  \code{member} has both efficiency and robustness advantages.
  \xlref{advanced-type-stuff}.
  
\item Type inference eliminates the need for some declarations, and
  also aids compile-time detection of type errors.  Given detailed
  type declarations, type inference can often eliminate type checks
  and enable more efficient object representations and code sequences.
  Checking all types results in fewer type checks.  See sections
  \ref{type-inference} and \ref{non-descriptor}.
\end{itemize}


\subsection{Optimization}

The main barrier to efficient Lisp programs is not that there is no
efficient way to code the program in Lisp, but that it is difficult to
arrive at that efficient coding.  \clisp{} is a highly complex
language, and usually has many semantically equivalent ``reasonable''
ways to code a given problem.  It is desirable to make all of these
equivalent solutions have comparable efficiency so that programmers
don't have to waste time discovering the most efficient solution.

Source level optimization increases the number of efficient ways to
solve a problem.  This effect is much larger than the increase in the
efficiency of the ``best'' solution.  Source level optimization
transforms the original program into a more efficient (but equivalent)
program.  Although the optimizer isn't doing anything the programmer
couldn't have done, this high-level optimization is important because:

\begin{itemize} 
\item The programmer can code simply and directly, rather than
  obfuscating code to please the compiler.
  
\item When presented with a choice of similar coding alternatives, the
  programmer can chose whichever happens to be most convenient,
  instead of worrying about which is most efficient.
\end{itemize}

Source level optimization eliminates the need for macros to optimize
their expansion, and also increases the effectiveness of inline
expansion.  See sections \ref{source-optimization} and
\ref{inline-expansion}.

Efficient support for a safer programming style is the biggest
advantage of source level optimization.  Existing tuned programs
typically won't benefit much from source optimization, since their
source has already been optimized by hand.  However, even tuned
programs tend to run faster under \python{} because:

\begin{itemize} 
\item Low level optimization and register allocation provides modest
  speedups in any program.
  
\item Block compilation and inline expansion can reduce function call
  overhead, but may require some program restructuring.  See sections
  \ref{inline-expansion}, \ref{local-call} and
  \ref{block-compilation}.
  
\item Efficiency notes will point out important type declarations that
  are often missed even in highly tuned programs.
  \xlref{efficiency-notes}.
  
\item Existing programs can be compiled safely without prohibitive
  speed penalty, although they would be faster and safer with added
  declarations.  \xlref{type-check-optimization}.
  
\item The context declaration mechanism allows both space and runtime
  of large systems to be reduced without sacrificing robustness by
  semi-automatically varying compilation policy without addition any
  \code{optimize} declarations to the source.
  \xlref{context-declarations}.
  
\item Byte compilation can be used to dramatically reduce the size of
  code that is not speed-critical. \xlref{byte-compile}
\end{itemize}


\subsection{Function Call}

The sort of symbolic programs generally written in \llisp{} often
favor recursion over iteration, or have inner loops so complex that
they involve multiple function calls.  Such programs spend a larger
fraction of their time doing function calls than is the norm in other
languages; for this reason \llisp{} implementations strive to make the
general (or full) function call as inexpensive as possible.  \python{}
goes beyond this by providing two good alternatives to full call:

\begin{itemize} 
\item Local call resolves function references at compile time,
  allowing better calling sequences and optimization across function
  calls.  \xlref{local-call}.
  
\item Inline expansion totally eliminates call overhead and allows
  many context dependent optimizations.  This provides a safe and
  efficient implementation of operations with function semantics,
  eliminating the need for error-prone macro definitions or manual
  case analysis.  Although most \clisp{} implementations support
  inline expansion, it becomes a more powerful tool with \python{}'s
  source level optimization.  See sections \ref{source-optimization}
  and \ref{inline-expansion}.
\end{itemize}


Generally, \python{} provides simple implementations for simple uses
of function call, rather than having only a single calling convention.
These features allow a more natural programming style:

\begin{itemize} 
\item Proper tail recursion.  \xlref{tail-recursion}
  
\item Relatively efficient closures.
  
\item A \code{funcall} that is as efficient as normal named call.
  
\item Calls to local functions such as from \code{labels} are
  optimized:
\begin{itemize}
  
\item Control transfer is a direct jump.
  
\item The closure environment is passed in registers rather than heap
  allocated.
  
\item Keyword arguments and multiple values are implemented more
  efficiently.
\end{itemize}

\xlref{local-call}.
\end{itemize}


\subsection{Representation of Objects}

Sometimes traditional \llisp{} implementation techniques compare so
poorly to the techniques used in other languages that \llisp{} can
become an impractical language choice.  Terrible inefficiencies appear
in number-crunching programs, since \llisp{} numeric operations often
involve number-consing and generic arithmetic.  \python{} supports
efficient natural representations for numbers (and some other types),
and allows these efficient representations to be used in more
contexts.  \python{} also provides good efficiency notes that warn
when a crucial declaration is missing.

See section \ref{non-descriptor} for more about object representations and
numeric types.  Also \pxlref{efficiency-notes} about efficiency notes.


\subsection{Writing Efficient Code}
\label{efficiency-overview}

Writing efficient code that works is a complex and prolonged process.
It is important not to get so involved in the pursuit of efficiency
that you lose sight of what the original problem demands.  Remember
that:
\begin{itemize}
  
\item The program should be correct\dash{}it doesn't matter how
  quickly you get the wrong answer.
  
\item Both the programmer and the user will make errors, so the
  program must be robust\dash{}it must detect errors in a way that
  allows easy correction.
  
\item A small portion of the program will consume most of the
  resources, with the bulk of the code being virtually irrelevant to
  efficiency considerations.  Even experienced programmers familiar
  with the problem area cannot reliably predict where these ``hot
  spots'' will be.
\end{itemize}



The best way to get efficient code that is still worth using, is to separate
coding from tuning.  During coding, you should:
\begin{itemize}
  
\item Use a coding style that aids correctness and robustness without
  being incompatible with efficiency.
  
\item Choose appropriate data structures that allow efficient
  algorithms and object representations
  (\pxlref{object-representation}).  Try to make interfaces abstract
  enough so that you can change to a different representation if
  profiling reveals a need.
  
\item Whenever you make an assumption about a function argument or
  global data structure, add consistency assertions, either with type
  declarations or explicit uses of \code{assert}, \code{ecase}, etc.
\end{itemize}

During tuning, you should:
\begin{itemize}
  
\item Identify the hot spots in the program through profiling (section
  \ref{profiling}.)
  
\item Identify inefficient constructs in the hot spot with efficiency
  notes, more profiling, or manual inspection of the source.  See
  sections \ref{general-efficiency} and \ref{efficiency-notes}.
  
\item Add declarations and consider the application of optimizations.
  See sections \ref{local-call}, \ref{inline-expansion} and
  \ref{non-descriptor}.
  
\item If all else fails, consider algorithm or data structure changes.
  If you did a good job coding, changes will be easy to introduce.
\end{itemize}


%%%% SBCL doesn't have a byte code interpreter. It's not to be ported
%%%% to the SBCL user manual. Commented out. -- MG

%% \section{Byte Coded Compilation}
%% \label{byte-compile}
%% \cindex{byte coded compilation}
%% \cindex{space optimization}
%% 
%% \python{} supports byte compilation to reduce the size of Lisp
%% programs by allowing functions to be compiled more compactly.  Byte
%% compilation provides an extreme speed/space tradeoff: byte code is
%% typically six times more compact than native code, but runs fifty
%% times (or more) slower.  This is about ten times faster than the
%% standard interpreter, which is itself considered fast in comparison to
%% other \clisp{} interpreters.
%% 
%% Large Lisp systems (such as \cmucl{} itself) often have large amounts
%% of user-interface code, compile-time (macro) code, debugging code, or
%% rarely executed special-case code.  This code is a good target for
%% byte compilation: very little time is spent running in it, but it can
%% take up quite a bit of space.  Straight-line code with many function
%% calls is much more suitable than inner loops.
%% 
%% When byte-compiling, the compiler compiles about twice as fast, and
%% can produce a hardware independent object file (\file{.bytef} type.)
%% This file can be loaded like a normal fasl file on any implementation
%% of \cmucl{} with the same byte-ordering.
%% 
%% The decision to byte compile or native compile can be done on a
%% per-file or per-code-object basis.  The \kwd{byte-compile} argument to
%% \funref{compile-file} has these possible values:
%% 
%% \begin{Lentry}
%% \item[\false{}] Don't byte compile anything in this file.
%%   
%% \item[\true{}] Byte compile everything in this file and produce a
%%   processor-independent \file{.bytef} file.
%%   
%% \item[\kwd{maybe}] Produce a normal fasl file, but byte compile any
%%   functions for which the \code{speed} optimization quality is
%%   \code{0} and the \code{debug} quality is not greater than \code{1}.
%% \end{Lentry}
%% 
%% \begin{defvar}{extensions:}{byte-compile-top-level}
%%   
%%   If this variable is true (the default) and the \kwd{byte-compile}
%%   argument to \code{compile-file} is \kwd{maybe}, then byte compile
%%   top-level code (code outside of any \code{defun}, \code{defmethod},
%%   etc.)
%% \end{defvar}
%% 
%% \begin{defvar}{extensions:}{byte-compile-default}
%%   
%%   This variable determines the default value for the
%%   \kwd{byte-compile} argument to \code{compile-file}, initially
%%   \kwd{maybe}.
%% \end{defvar}


%%%% This section is more like a Common Lisp tutorial. It's not to be
%%%% ported to the SBCL user manual. Commented out. -- MG

%% \section{Object Representation}
%% \label{object-representation}
%% \cindex{object representation}
%% \cpsubindex{representation}{object}
%% \cpsubindex{efficiency}{of objects}
%% 
%% A somewhat subtle aspect of writing efficient \clisp{} programs is
%% choosing the correct data structures so that the underlying objects
%% can be implemented efficiently.  This is partly because of the need
%% for multiple representations for a given value
%% (\pxlref{non-descriptor}), but is also due to the sheer number of
%% object types that \clisp{} has built in.  The number of possible
%% representations complicates the choice of a good representation
%% because semantically similar objects may vary in their efficiency
%% depending on how the program operates on them.
%% 
%% 
%% \subsection{Think Before You Use a List}
%% \cpsubindex{lists}{efficiency of}
%% 
%% Although Lisp's creator seemed to think that it was for LISt
%% Processing, the astute observer may have noticed that the chapter on
%% list manipulation makes up less that three percent of \cltltwo{}. The
%% language has grown since Lisp 1.5\dash{}new data types supersede lists
%% for many purposes.
%% 
%% 
%% \subsection{Structure Representation}
%% \cpsubindex{structure types}{efficiency of} One of the best ways of
%% building complex data structures is to define appropriate structure
%% types using \findexed{defstruct}.  In \python, access of structure
%% slots is always at least as fast as list or vector access, and is
%% usually faster.  In comparison to a list representation of a tuple,
%% structures also have a space advantage.
%% 
%% Even if structures weren't more efficient than other representations, structure
%% use would still be attractive because programs that use structures in
%% appropriate ways are much more maintainable and robust than programs written
%% using only lists.  For example:
%% \begin{lisp}
%% (rplaca (caddr (cadddr x)) (caddr y))
%% \end{lisp}
%% could have been written using structures in this way:
%% \begin{lisp}
%% (setf (beverage-flavor (astronaut-beverage x)) (beverage-flavor y))
%% \end{lisp}
%% The second version is more maintainable because it is easier to
%% understand what it is doing.  It is more robust because structures
%% accesses are type checked.  An \code{astronaut} will never be confused
%% with a \code{beverage}, and the result of \code{beverage-flavor} is
%% always a flavor.  See sections \ref{structure-types} and
%% \ref{freeze-type} for more information about structure types.
%% \xlref{type-inference} for a number of examples that make clear the
%% advantages of structure typing.
%% 
%% Note that the structure definition should be compiled before any uses
%% of its accessors or type predicate so that these function calls can be
%% efficiently open-coded.
%% 
%% 
%% \subsection{Arrays}
%% \label{array-types}
%% \cpsubindex{arrays}{efficiency of}
%% 
%% Arrays are often the most efficient representation for collections of objects
%% because:
%% \begin{itemize}
%%   
%% \item Array representations are often the most compact.  An array is
%%   always more compact than a list containing the same number of
%%   elements.
%%   
%% \item Arrays allow fast constant-time access.
%%   
%% \item Arrays are easily destructively modified, which can reduce
%%   consing.
%%   
%% \item Array element types can be specialized, which reduces both
%%   overall size and consing (\pxlref{specialized-array-types}.)
%% \end{itemize}
%% 
%% 
%% Access of arrays that are not of type \code{simple-array} is less
%% efficient, so declarations are appropriate when an array is of a
%% simple type like \code{simple-string} or \code{simple-bit-vector}.
%% Arrays are almost always simple, but the compiler may not be able to
%% prove simpleness at every use.  The only way to get a non-simple array
%% is to use the \kwd{displaced-to}, \kwd{fill-pointer} or
%% \code{adjustable} arguments to \code{make-array}.  If you don't use
%% these hairy options, then arrays can always be declared to be simple.
%% 
%% Because of the many specialized array types and the possibility of
%% non-simple arrays, array access is much like generic arithmetic
%% (\pxlref{generic-arithmetic}).  In order for array accesses to be
%% efficiently compiled, the element type and simpleness of the array
%% must be known at compile time.  If there is inadequate information,
%% the compiler is forced to call a generic array access routine.  You
%% can detect inefficient array accesses by enabling efficiency notes,
%% \pxlref{efficiency-notes}.
%% 
%% 
%% \subsection{Vectors}
%% \cpsubindex{vectors}{efficiency of}
%% 
%% Vectors (one dimensional arrays) are particularly useful, since in
%% addition to their obvious array-like applications, they are also well
%% suited to representing sequences.  In comparison to a list
%% representation, vectors are faster to access and take up between two
%% and sixty-four times less space (depending on the element type.)  As
%% with arbitrary arrays, the compiler needs to know that vectors are not
%% complex, so you should use \code{simple-string} in preference to
%% \code{string}, etc.
%% 
%% The only advantage that lists have over vectors for representing
%% sequences is that it is easy to change the length of a list, add to it
%% and remove items from it.  Likely signs of archaic, slow lisp code are
%% \code{nth} and \code{nthcdr}.  If you are using these functions you
%% should probably be using a vector.
%% 
%% 
%% \subsection{Bit-Vectors}
%% \cpsubindex{bit-vectors}{efficiency of}
%% 
%% Another thing that lists have been used for is set manipulation.  In
%% applications where there is a known, reasonably small universe of
%% items bit-vectors can be used to improve performance.  This is much
%% less convenient than using lists, because instead of symbols, each
%% element in the universe must be assigned a numeric index into the bit
%% vector.  Using a bit-vector will nearly always be faster, and can be
%% tremendously faster if the number of elements in the set is not small.
%% The logical operations on \code{simple-bit-vector}s are efficient,
%% since they operate on a word at a time.
%% 
%% 
%% \subsection{Hashtables}
%% \cpsubindex{hash-tables}{efficiency of}
%% 
%% Hashtables are an efficient and general mechanism for maintaining associations
%% such as the association between an object and its name.  Although hashtables
%% are usually the best way to maintain associations, efficiency and style
%% considerations sometimes favor the use of an association list (a-list).
%% 
%% \code{assoc} is fairly fast when the \var{test} argument is \code{eq}
%% or \code{eql} and there are only a few elements, but the time goes up
%% in proportion with the number of elements.  In contrast, the
%% hash-table lookup has a somewhat higher overhead, but the speed is
%% largely unaffected by the number of entries in the table.  For an
%% \code{equal} hash-table or alist, hash-tables have an even greater
%% advantage, since the test is more expensive.  Whatever you do, be sure
%% to use the most restrictive test function possible.
%% 
%% The style argument observes that although hash-tables and alists
%% overlap in function, they do not do all things equally well.
%% \begin{itemize}
%%   
%% \item Alists are good for maintaining scoped environments.  They were
%%   originally invented to implement scoping in the Lisp interpreter,
%%   and are still used for this in \python.  With an alist one can
%%   non-destructively change an association simply by consing a new
%%   element on the front.  This is something that cannot be done with
%%   hash-tables.
%%   
%% \item Hashtables are good for maintaining a global association.  The
%%   value associated with an entry can easily be changed with
%%   \code{setf}.  With an alist, one has to go through contortions,
%%   either \code{rplacd}'ing the cons if the entry exists, or pushing a
%%   new one if it doesn't.  The side-effecting nature of hash-table
%%   operations is an advantage here.
%% \end{itemize}
%% 
%% 
%% Historically, symbol property lists were often used for global name
%% associations.  Property lists provide an awkward and error-prone
%% combination of name association and record structure.  If you must use
%% the property list, please store all the related values in a single
%% structure under a single property, rather than using many properties.
%% This makes access more efficient, and also adds a modicum of typing
%% and abstraction.  \xlref{advanced-type-stuff} for information on types
%% in \cmucl.


\section{Numbers}
\label{numeric-types}
\cpsubindex{numeric}{types}
\cpsubindex{types}{numeric}

Numbers are interesting because numbers are one of the few \llisp{} data types
that have direct support in conventional hardware.  If a number can be
represented in the way that the hardware expects it, then there is a big
efficiency advantage.

Using hardware representations is problematical in \llisp{} due to
dynamic typing (where the type of a value may be unknown at compile
time.)  It is possible to compile code for statically typed portions
of a \llisp{} program with efficiency comparable to that obtained in
statically typed languages such as C, but not all \llisp{}
implementations succeed.  There are two main barriers to efficient
numerical code in \llisp{}:
\begin{itemize}
  
\item The compiler must prove that the numerical expression is in fact
  statically typed, and
  
\item The compiler must be able to somehow reconcile the conflicting
  demands of the hardware mandated number representation with the
  \llisp{} requirements of dynamic typing and garbage-collecting
  dynamic storage allocation.
\end{itemize}

Because of its type inference (\pxlref{type-inference}) and efficiency
notes (\pxlref{efficiency-notes}), \python{} is better than
conventional \llisp{} compilers at ensuring that numerical expressions
are statically typed.  \python{} also goes somewhat farther than existing
compilers in the area of allowing native machine number
representations in the presence of garbage collection.


\subsection{Descriptors}
\cpsubindex{descriptors}{object}
\cindex{object representation}
\cpsubindex{representation}{object}
\cpsubindex{consing}{overhead of}

\llisp{}'s dynamic typing requires that it be possible to represent
any value with a fixed length object, known as a \var{descriptor}.
This fixed-length requirement is implicit in features such as:
\begin{itemize}
  
\item Data types (like \code{simple-vector}) that can contain any type
  of object, and that can be destructively modified to contain
  different objects (of possibly different types.)
  
\item Functions that can be called with any type of argument, and that
  can be redefined at run time.
\end{itemize}

In order to save space, a descriptor is invariably represented as a
single word.  Objects that can be directly represented in the
descriptor itself are said to be \var{immediate}.  Descriptors for
objects larger than one word are in reality pointers to the memory
actually containing the object.

Representing objects using pointers has two major disadvantages:
\begin{itemize}
  
\item The memory pointed to must be allocated on the heap, so it must
  eventually be freed by the garbage collector.  Excessive heap
  allocation of objects (or ``consing'') is inefficient in several
  ways.  \xlref{consing}.
  
\item Representing an object in memory requires the compiler to emit
  additional instructions to read the actual value in from memory, and
  then to write the value back after operating on it.
\end{itemize}

The introduction of garbage collection makes things even worse, since
the garbage collector must be able to determine whether a descriptor
is an immediate object or a pointer.  This requires that a few bits in
each descriptor be dedicated to the garbage collector.  The loss of a
few bits doesn't seem like much, but it has a major efficiency
implication\dash{}objects whose natural machine representation is a
full word (integers and single-floats) cannot have an immediate
representation.  So the compiler is forced to use an unnatural
immediate representation (such as \code{fixnum}) or a natural pointer
representation (with the attendant consing overhead.)


\subsection{Non-Descriptor Representations}
\label{non-descriptor}
\cindex{non-descriptor representations}
\cindex{stack numbers}

From the discussion above, we can see that the standard descriptor
representation has many problems, the worst being number consing.
\llisp{} compilers try to avoid these descriptor efficiency problems by using
\var{non-descriptor} representations.  A compiler that uses non-descriptor
representations can compile this function so that it does no number consing:
\begin{lisp}
(defun multby (vec n)
  (declare (type (simple-array single-float (*)) vec)
           (single-float n))
  (dotimes (i (length vec))
    (setf (aref vec i)
          (* n (aref vec i)))))
\end{lisp}
If a descriptor representation were used, each iteration of the loop might
cons two floats and do three times as many memory references.

As its negative definition suggests, the range of possible non-descriptor
representations is large.  The performance improvement from non-descriptor
representation depends upon both the number of types that have non-descriptor
representations and the number of contexts in which the compiler is forced to
use a descriptor representation.

Many \llisp{} compilers support non-descriptor representations for
float types such as \code{single-float} and \code{double-float}
(section \ref{float-efficiency}.)  \python{} adds support for full
word integers (\pxlref{word-integers}), characters
(\pxlref{characters}) and system-area pointers (unconstrained
pointers, \pxlref{system-area-pointers}.)  Many \llisp{} compilers
support non-descriptor representations for variables (section
\ref{ND-variables}) and array elements (section
\ref{specialized-array-types}.)  \python{} adds support for
non-descriptor arguments and return values in local call
(\pxlref{number-local-call}) and structure slots (\pxlref{raw-slots}).


\subsection{Variables}
\label{ND-variables}
\cpsubindex{variables}{non-descriptor}
\cpsubindex{type declarations}{variable}
\cpsubindex{efficiency}{of numeric variables}

In order to use a non-descriptor representation for a variable or
expression intermediate value, the compiler must be able to prove that
the value is always of a particular type having a non-descriptor
representation.  Type inference (\pxlref{type-inference}) often needs
some help from user-supplied declarations.  The best kind of type
declaration is a variable type declaration placed at the binding
point:
\begin{lisp}
(let ((x (car l)))
  (declare (single-float x))
  ...)
\end{lisp}
Use of \code{the}, or of variable declarations not at the binding form
is insufficient to allow non-descriptor representation of the
variable\dash{}with these declarations it is not certain that all
values of the variable are of the right type.  It is sometimes useful
to introduce a gratuitous binding that allows the compiler to change
to a non-descriptor representation, like:
\begin{lisp}
(etypecase x
  ((signed-byte 32)
   (let ((x x))
     (declare (type (signed-byte 32) x)) 
     ...))
  ...)
\end{lisp}
The declaration on the inner \code{x} is necessary here due to a phase
ordering problem.  Although the compiler will eventually prove that
the outer \code{x} is a \w{\code{(signed-byte 32)}} within that
\code{etypecase} branch, the inner \code{x} would have been optimized
away by that time.  Declaring the type makes let optimization more
cautious.

Note that storing a value into a global (or \code{special}) variable
always forces a descriptor representation.  Wherever possible, you
should operate only on local variables, binding any referenced globals
to local variables at the beginning of the function, and doing any
global assignments at the end.

Efficiency notes signal use of inefficient representations, so
programmer's needn't continuously worry about the details of
representation selection (\pxlref{representation-eff-note}.)


\subsection{Generic Arithmetic}
\label{generic-arithmetic}
\cindex{generic arithmetic}
\cpsubindex{arithmetic}{generic}
\cpsubindex{numeric}{operation efficiency}

In \clisp, arithmetic operations are \var{generic}.\footnote{As Steele
  notes in CLTL II, this is a generic conception of generic, and is
  not to be confused with the CLOS concept of a generic function.}
The \code{+} function can be passed \code{fixnum}s, \code{bignum}s,
\code{ratio}s, and various kinds of \code{float}s and
\code{complex}es, in any combination.  In addition to the inherent
complexity of \code{bignum} and \code{ratio} operations, there is also
a lot of overhead in just figuring out which operation to do and what
contagion and canonicalization rules apply.  The complexity of generic
arithmetic is so great that it is inconceivable to open code it.
Instead, the compiler does a function call to a generic arithmetic
routine, consuming many instructions before the actual computation
even starts.

This is ridiculous, since even \llisp{} programs do a lot of
arithmetic, and the hardware is capable of doing operations on small
integers and floats with a single instruction.  To get acceptable
efficiency, the compiler special-cases uses of generic arithmetic that
are directly implemented in the hardware.  In order to open code
arithmetic, several constraints must be met:
\begin{itemize}
  
\item All the arguments must be known to be a good type of number.
  
\item The result must be known to be a good type of number.
  
\item Any intermediate values such as the result of \w{\code{(+ a b)}}
  in the call \w{\code{(+ a b c)}} must be known to be a good type of
  number.
  
\item All the above numbers with good types must be of the \var{same}
  good type.  Don't try to mix integers and floats or different float
  formats.
\end{itemize}

The ``good types'' are \w{\code{(signed-byte 32)}},
\w{\code{(unsigned-byte 32)}}, \code{single-float},
\code{double-float}, \code{(complex single-float)}, and \code{(complex
  double-float)}.  See sections \ref{fixnums}, \ref{word-integers}
and \ref{float-efficiency} for more discussion of good numeric types.

\code{float} is not a good type, since it might mean either
\code{single-float} or \code{double-float}.  \code{integer} is not a
good type, since it might mean \code{bignum}.  \code{rational} is not
a good type, since it might mean \code{ratio}.  Note however that
these types are still useful in declarations, since type inference may
be able to strengthen a weak declaration into a good one, when it
would be at a loss if there was no declaration at all
(\pxlref{type-inference}).  The \code{integer} and
\code{unsigned-byte} (or non-negative integer) types are especially
useful in this regard, since they can often be strengthened to a good
integer type.

% Arithmetic with \code{complex} numbers is inefficient in comparison to
% float and integer arithmetic.  Complex numbers are always represented
% with a pointer descriptor (causing consing overhead), and complex
% arithmetic is always closed coded using the general generic arithmetic
As noted above, \cmucl{} has support for \code{(complex single-float)}
and \code{(complex double-float)}.  These can be unboxed and, thus,
are quite efficient.  However, arithmetic with complex types such as:
\begin{lisp}
(complex float)
(complex fixnum)
\end{lisp}
will be significantly slower than the good complex types but is still
faster than \code{bignum} or \code{ratio} arithmetic, since the
implementation is much simpler.

Note: don't use \code{/} to divide integers unless you want the
overhead of rational arithmetic.  Use \code{truncate} even when you
know that the arguments divide evenly.

You don't need to remember all the rules for how to get open-coded
arithmetic, since efficiency notes will tell you when and where there
is a problem\dash{}\pxlref{efficiency-notes}.


\subsection{Fixnums}
\label{fixnums}
\cindex{fixnums}
\cindex{bignums}

A fixnum is a ``FIXed precision NUMber''.  In modern \llisp{}
implementations, fixnums can be represented with an immediate
descriptor, so operating on fixnums requires no consing or memory
references.  Clever choice of representations also allows some
arithmetic operations to be done on fixnums using hardware supported
word-integer instructions, somewhat reducing the speed penalty for
using an unnatural integer representation.

It is useful to distinguish the \code{fixnum} type from the fixnum
representation of integers.  In \python, there is absolutely nothing
magical about the \code{fixnum} type in comparison to other finite
integer types.  \code{fixnum} is equivalent to (is defined with
\code{deftype} to be) \w{\code{(signed-byte 30)}}.  \code{fixnum} is
simply the largest subset of integers that {\em can be represented}
using an immediate fixnum descriptor.

Unlike in other \clisp{} compilers, it is in no way desirable to use
the \code{fixnum} type in declarations in preference to more
restrictive integer types such as \code{bit}, \w{\code{(integer -43
    7)}} and \w{\code{(unsigned-byte 8)}}.  Since \python{} does
understand these integer types, it is preferable to use the more
restrictive type, as it allows better type inference
(\pxlref{operation-type-inference}.)

The small, efficient fixnum is contrasted with bignum, or ``BIG
NUMber''.  This is another descriptor representation for integers, but
this time a pointer representation that allows for arbitrarily large
integers.  Bignum operations are less efficient than fixnum
operations, both because of the consing and memory reference overheads
of a pointer descriptor, and also because of the inherent complexity
of extended precision arithmetic.  While fixnum operations can often
be done with a single instruction, bignum operations are so complex
that they are always done using generic arithmetic.

A crucial point is that the compiler will use generic arithmetic if it
can't \var{prove} that all the arguments, intermediate values, and
results are fixnums.  With bounded integer types such as
\code{fixnum}, the result type proves to be especially problematical,
since these types are not closed under common arithmetic operations
such as \code{+}, \code{-}, \code{*} and \code{/}.  For example,
\w{\code{(1+ (the fixnum x))}} does not necessarily evaluate to a
\code{fixnum}.  Bignums were added to \llisp{} to get around this
problem, but they really just transform the correctness problem ``if
this add overflows, you will get the wrong answer'' to the efficiency
problem ``if this add \var{might} overflow then your program will run
slowly (because of generic arithmetic.)''

There is just no getting around the fact that the hardware only
directly supports short integers.  To get the most efficient open
coding, the compiler must be able to prove that the result is a good
integer type.  This is an argument in favor of using more restrictive
integer types: \w{\code{(1+ (the fixnum x))}} may not always be a
\code{fixnum}, but \w{\code{(1+ (the (unsigned-byte 8) x))}} always
is.  Of course, you can also assert the result type by putting in lots
of \code{the} declarations and then compiling with \code{safety}
\code{0}.


\subsection{Word Integers}
\label{word-integers}
\cindex{word integers}

\python{} is unique in its efficient implementation of arithmetic
on full-word integers through non-descriptor representations and open coding.
Arithmetic on any subtype of these types:

\begin{lisp}
(signed-byte 32)
(unsigned-byte 32)
\end{lisp}

is reasonably efficient, although subtypes of \code{fixnum} remain
somewhat more efficient.

If a word integer must be represented as a descriptor, then the
\code{bignum} representation is used, with its associated consing
overhead.  The support for word integers in no way changes the
language semantics, it just makes arithmetic on small bignums vastly
more efficient.  It is fine to do arithmetic operations with mixed
\code{fixnum} and word integer operands; just declare the most
specific integer type you can, and let the compiler decide what
representation to use.

In fact, to most users, the greatest advantage of word integer
arithmetic is that it effectively provides a few guard bits on the
fixnum representation.  If there are missing assertions on
intermediate values in a fixnum expression, the intermediate results
can usually be proved to fit in a word.  After the whole expression is
evaluated, there will often be a fixnum assertion on the final result,
allowing creation of a fixnum result without even checking for
overflow.

The remarks in section \ref{fixnums} about fixnum result type also
apply to word integers; you must be careful to give the compiler
enough information to prove that the result is still a word integer.
This time, though, when we blow out of word integers we land in into
generic bignum arithmetic, which is much worse than sleazing from
\code{fixnum}s to word integers.  Note that mixing
\w{\code{(unsigned-byte 32)}} arguments with arguments of any signed
type (such as \code{fixnum}) is a no-no, since the result might not be
unsigned.


\subsection{Floating Point Efficiency}
\label{float-efficiency}
\cindex{floating point efficiency}

Arithmetic on objects of type \code{single-float} and \code{double-float} is
efficiently implemented using non-descriptor representations and open coding.
As for integer arithmetic, the arguments must be known to be of the same float
type.  Unlike for integer arithmetic, the results and intermediate values
usually take care of themselves due to the rules of float contagion, i.e.
\w{\code{(1+ (the single-float x))}} is always a \code{single-float}.

Although they are not specially implemented, \code{short-float} and
\code{long-float} are also acceptable in declarations, since they are
synonyms for the \code{single-float} and \code{double-float} types,
respectively.

In \cmucl{}, list-style float type specifiers such as
\w{\code{(single-float 0.0 1.0)}} will be used to good effect.

For example, in this function,

\begin{example}
  (defun square (x)
    (declare (type (single-float 0f0 10f0)))
    (* x x))
\end{example}

\python{} can deduce that the
return type of the function \code{square} is \w{\code{(single-float
    0f0 100f0)}}.

Many union types are also supported so that

\begin{example}
  (+ (the (or (integer 1 1) (integer 5 5)) x)
     (the (or (integer 10 10) (integer 20 20)) y))
\end{example}

has the inferred type \code{(or (integer 11 11) (integer 15 15)
  (integer 21 21) (integer 25 25))}.  This also works for
floating-point numbers.  Member types are also supported.
  
\cmucl{} can also infer types for many mathematical functions
including square root, exponential and logarithmic functions,
trignometric functions and their inverses, and hyperbolic functions
and their inverses.  For numeric code, this can greatly enhance
efficiency by allowing the compiler to use specialized versions of
the functions instead of the generic versions.  The greatest benefit 
of this type inference is determining that the result of the
function is real-valued number instead of possibly being
a complex-valued number.

For example, consider the function
\begin{example}
  (defun fun (x)
    (declare (type (single-float (0f0) 100f0) x))
    (values (sqrt x) (log x)))
\end{example}
With this declaration, the compiler can determine that the argument
to \code{sqrt} and \code{log} are always non-negative so that the result 
is always a \code{single-float}.  In fact, the return type for this
function is derived to be \code{(values (single-float 0f0 10f0)
    (single-float * 2f0))}.

If the declaration were reduced to just \w{\code{(declare
    (single-float x))}}, the argument to \code{sqrt} and \code{log}
could be negative.  This forces the use of the generic versions of
these functions because the result could be a complex number.

We note, however, that proper interval arithmetic is not fully
implemented in the compiler so the inferred types may be slightly in
error due to round-off errors.  This round-off error could
accumulate to cause the compiler to erroneously deduce the result
type and cause code to be removed as being
unreachable.\footnote{This, however, has not actually happened, but
  it is a possibility.}%
Thus, the declarations should only be precise enough for the
compiler to deduce that a real-valued argument to a function would
produce a real-valued result.  The efficiency notes
(\pxlref{representation-eff-note}) from the compiler will guide you
on what declarations might be useful.

When a float must be represented as a descriptor, a pointer representation is
used, creating consing overhead.  For this reason, you should try to avoid
situations (such as full call and non-specialized data structures) that force a
descriptor representation.  See sections \ref{specialized-array-types},
\ref{raw-slots} and \ref{number-local-call}.

\xlref{ieee-float} for information on the extensions to support IEEE
floating point.

\subsubsection{Signed Zeroes and Special Functions}

\cmucl{} supports IEEE signed zeroes.  In typical usage, the signed
zeroes are not a problem and can be treated as an unsigned zero.
However, some of the special functions have branch points at zero, so
care must be taken.

For example, suppose we have the function
\begin{example}
  (defun fun (x)
    (declare (type (single-float 0f0) x))
    (log x))
\end{example}
The derived result of the function is \code{(OR SINGLE-FLOAT
(COMPLEX SINGLE-FLOAT))} because the declared values for
\code{x} includes both $-0.0$ and $0.0$ and \code{(log -0.0)} is
actually a complex number.  Because of this, the generic complex log
routine is used.

If the declaration for \code{x} were \code{(single-float (0f0))} so $+0.0$
is not included or \code{(or (single-float (0f0)) (member 0f0))} so
  $+0.0$ is include but not $-0.0$, the derived type would be
  \code{single-float} for both cases.  By declaring \code{x} this way,
  the log can be implemented using a fast real-valued log routine
  instead of the generic log routine.

\cmucl{} implements the branch cuts and values given by
Kahan\footnote{Kahan, W., ``Branch Cuts for Complex Elementary
Functions, or Much Ado About Nothing's Sign Bit''
in Iserles and Powell (eds.) \textit{The State of the Art
in Numerical Analysis}, pp. 165-211, Clarendon
Press, 1987}.
  


\subsection{Specialized Arrays}
\label{specialized-array-types}
\cindex{specialized array types}
\cpsubindex{array types}{specialized}
\cpsubindex{types}{specialized array}

\clisp{} supports specialized array element types through the
\kwd{element-type} argument to \code{make-array}.  When an array has a
specialized element type, only elements of that type can be stored in
the array.  From this restriction comes two major efficiency
advantages:
\begin{itemize}
  
\item A specialized array can save space by packing multiple elements
  into a single word.  For example, a \code{base-char} array can have
  4 elements per word, and a \code{bit} array can have 32.  This
  space-efficient representation is possible because it is not
  necessary to separately indicate the type of each element.
  
\item The elements in a specialized array can be given the same
  non-descriptor representation as the one used in registers and on
  the stack, eliminating the need for representation conversions when
  reading and writing array elements.  For objects with pointer
  descriptor representations (such as floats and word integers) there
  is also a substantial consing reduction because it is not necessary
  to allocate a new object every time an array element is modified.
\end{itemize}


These are the specialized element types currently supported:
\begin{lisp}
bit
(unsigned-byte 2)
(unsigned-byte 4)
(unsigned-byte 8)
(unsigned-byte 16)
(unsigned-byte 32)
(signed-byte 8)
(signed-byte 16)
(signed-byte 30)
(signed-byte 32)
base-character
single-float
double-float
ext:double-double-float
(complex single-float)
(complex double-float)
(complex ext:double-double-float)
\end{lisp}

Although a \code{simple-vector} can hold any type of object, \true{}
should still be considered a specialized array type, since arrays with
element type \true{} are specialized to hold descriptors.



When using non-descriptor representations, it is particularly
important to make sure that array accesses are open-coded, since in
addition to the generic operation overhead, efficiency is lost when
the array element is converted to a descriptor so that it can be
passed to (or from) the generic access routine.  You can detect
inefficient array accesses by enabling efficiency notes,
\pxlref{efficiency-notes}.  \xlref{array-types}.


\subsection{Specialized Structure Slots}
\label{raw-slots}
\cpsubindex{structure types}{numeric slots}
\cindex{specialized structure slots}

Structure slots declared by the \kwd{type} \code{defstruct} slot option
to have certain known numeric types are also given non-descriptor
representations.  These types (and subtypes of these types) are supported:
\begin{lisp}
(unsigned-byte 32)
single-float
double-float
(complex single-float)
(complex double-float)
\end{lisp}

The primary advantage of specialized slot representations is a large
reduction spurious memory allocation and access overhead of programs
that intensively use these types.


\subsection{Interactions With Local Call}
\label{number-local-call}
\cpsubindex{local call}{numeric operands}
\cpsubindex{call}{numeric operands}
\cindex{numbers in local call}

Local call has many advantages (\pxlref{local-call}); one relevant to
our discussion here is that local call extends the usefulness of
non-descriptor representations.  If the compiler knows from the
argument type that an argument has a non-descriptor representation,
then the argument will be passed in that representation.  The easiest
way to ensure that the argument type is known at compile time is to
always declare the argument type in the called function, like:
\begin{lisp}
(defun 2+f (x)
  (declare (single-float x))
  (+ x 2.0))
\end{lisp}
The advantages of passing arguments and return values in a non-descriptor
representation are the same as for non-descriptor representations in general:
reduced consing and memory access (\pxlref{non-descriptor}.)  This
extends the applicative programming styles discussed in section
\ref{local-call} to numeric code.  Also, if source files are kept reasonably
small, block compilation can be used to reduce number consing to a minimum.

Note that non-descriptor return values can only be used with the known return
convention (section \ref{local-call-return}.)  If the compiler can't prove that
a function always returns the same number of values, then it must use the
unknown values return convention, which requires a descriptor representation.
Pay attention to the known return efficiency notes to avoid number consing.
 

\subsection{Representation of Characters}
\label{characters}
\cindex{characters}
\cindex{strings}

\python{} also uses a non-descriptor representation for characters when
convenient.  This improves the efficiency of string manipulation, but is
otherwise pretty invisible; characters have an immediate descriptor
representation, so there is not a great penalty for converting a character to a
descriptor.  Nonetheless, it may sometimes be helpful to declare
character-valued variables as \code{base-character}.


\section{General Efficiency Hints}
\label{general-efficiency}
\cpsubindex{efficiency}{general hints}

This section is a summary of various implementation costs and ways to get
around them.  These hints are relatively unrelated to the use of the \python{}
compiler, and probably also apply to most other \llisp{} implementations.  In
each section, there are references to related in-depth discussion.


\subsection{Compile Your Code}
\cpsubindex{compilation}{why to}

At this point, the advantages of compiling code relative to running it
interpreted probably need not be emphasized too much, but remember that
in \cmucl, compiled code typically runs hundreds of times faster than
interpreted code.  Also, compiled (\code{fasl}) files load significantly faster
than source files, so it is worthwhile compiling files which are loaded many
times, even if the speed of the functions in the file is unimportant.

Even disregarding the efficiency advantages, compiled code is as good or better
than interpreted code.  Compiled code can be debugged at the source level (see
chapter \ref{debugger}), and compiled code does more error checking.  For these
reasons, the interpreter should be regarded mainly as an interactive command
interpreter, rather than as a programming language implementation.

\b{Do not} be concerned about the performance of your program until you
see its speed compiled.  Some techniques that make compiled code run
faster make interpreted code run slower.


\subsection{Avoid Unnecessary Consing}
\label{consing}
\cindex{consing}
\cindex{garbage collection}
\cindex{memory allocation}
\cpsubindex{efficiency}{of memory use}


Consing is another name for allocation of storage, as done by the
\code{cons} function (hence its name.)  \code{cons} is by no means the
only function which conses\dash{}so does \code{make-array} and many
other functions.  Arithmetic and function call can also have hidden
consing overheads.  Consing hurts performance in the following ways:
\begin{itemize}
  
\item Consing reduces memory access locality, increasing paging
  activity.
  
\item Consing takes time just like anything else.
  
\item Any space allocated eventually needs to be reclaimed, either by
  garbage collection or by starting a new \code{lisp} process.
\end{itemize}


Consing is not undiluted evil, since programs do things other than
consing, and appropriate consing can speed up the real work.  It would
certainly save time to allocate a vector of intermediate results that
are reused hundreds of times.  Also, if it is necessary to copy a
large data structure many times, it may be more efficient to update
the data structure non-destructively; this somewhat increases update
overhead, but makes copying trivial.

Note that the remarks in section \ref{efficiency-overview} about the
importance of separating tuning from coding also apply to consing
overhead.  The majority of consing will be done by a small portion of
the program.  The consing hot spots are even less predictable than the
CPU hot spots, so don't waste time and create bugs by doing
unnecessary consing optimization.  During initial coding, avoid
unnecessary side-effects and cons where it is convenient.  If
profiling reveals a consing problem, \var{then} go back and fix the
hot spots.

\xlref{non-descriptor} for a discussion of how to avoid number consing
in \python.


\subsection{Complex Argument Syntax}
\cpsubindex{argument syntax}{efficiency}
\cpsubindex{efficiency}{of argument syntax}
\cindex{keyword argument efficiency}
\cindex{rest argument efficiency}

\clisp{} has very powerful argument passing mechanisms.  Unfortunately, two
of the most powerful mechanisms, rest arguments and keyword arguments, have a
significant performance penalty:

\begin{itemize}
\item
With keyword arguments, the called function has to parse the supplied keywords
by iterating over them and checking them against the desired keywords.

\item
With rest arguments, the function must cons a list to hold the arguments.  If a
function is called many times or with many arguments, large amounts of memory
will be allocated.
\end{itemize}

Although rest argument consing is worse than keyword parsing, neither problem
is serious unless thousands of calls are made to such a function.  The use of
keyword arguments is strongly encouraged in functions with many arguments or
with interfaces that are likely to be extended, and rest arguments are often
natural in user interface functions.

Optional arguments have some efficiency advantage over keyword
arguments, but their syntactic clumsiness and lack of extensibility
has caused many \clisp{} programmers to abandon use of optionals
except in functions that have obviously simple and immutable
interfaces (such as \code{subseq}), or in functions that are only
called in a few places.  When defining an interface function to be
used by other programmers or users, use of only required and keyword
arguments is recommended.

Parsing of \code{defmacro} keyword and rest arguments is done at
compile time, so a macro can be used to provide a convenient syntax
with an efficient implementation.  If the macro-expanded form contains
no keyword or rest arguments, then it is perfectly acceptable in inner
loops.

Keyword argument parsing overhead can also be avoided by use of inline
expansion (\pxlref{inline-expansion}) and block compilation (section
\ref{block-compilation}.)

Note: the compiler open-codes most heavily used system functions which have
keyword or rest arguments, so that no run-time overhead is involved.


\subsection{Mapping and Iteration}
\cpsubindex{mapping}{efficiency of}

One of the traditional \llisp{} programming styles is a highly applicative one,
involving the use of mapping functions and many lists to store intermediate
results.  To compute the sum of the square-roots of a list of numbers, one
might say:

\begin{lisp}
(apply #'+ (mapcar #'sqrt list-of-numbers))
\end{lisp}

This programming style is clear and elegant, but unfortunately results
in slow code.  There are two reasons why:

\begin{itemize} 
\item The creation of lists of intermediate results causes much
  consing (see \ref{consing}).
  
\item Each level of application requires another scan down the list.
  Thus, disregarding other effects, the above code would probably take
  twice as long as a straightforward iterative version.
\end{itemize}


An example of an iterative version of the same code:
\begin{lisp}
(do ((num list-of-numbers (cdr num))
     (sum 0 (+ (sqrt (car num)) sum)))
    ((null num) sum))
\end{lisp}

See sections \ref{variable-type-inference} and \ref{let-optimization}
for a discussion of the interactions of iteration constructs with type
inference and variable optimization.  Also, section
\ref{local-tail-recursion} discusses an applicative style of
iteration.


\subsection{Trace Files and Disassembly}
\label{trace-files}
\cindex{trace files}
\cindex{assembly listing}
\cpsubindex{listing files}{trace}
\cindex{Virtual Machine (VM, or IR2) representation}
\cindex{implicit continuation representation (IR1)}
\cpsubindex{continuations}{implicit representation}

In order to write efficient code, you need to know the relative costs
of different operations.  The main reason why writing efficient
\llisp{} code is difficult is that there are so many operations, and
the costs of these operations vary in obscure context-dependent ways.
Although efficiency notes point out some problem areas, the only way
to ensure generation of the best code is to look at the assembly code
output.

The \code{disassemble} function is a convenient way to get the assembly code for a
function, but it can be very difficult to interpret, since the correspondence
with the original source code is weak.  A better (but more awkward) option is
to use the \kwd{trace-file} argument to \code{compile-file} to generate a trace
file.

A trace file is a dump of the compiler's internal representations,
including annotated assembly code.  Each component in the program gets
four pages in the trace file (separated by ``\code{$\hat{ }L$}''):
\begin{itemize}
  
\item The implicit-continuation (or IR1) representation of the
  optimized source.  This is a dump of the flow graph representation
  used for ``source level'' optimizations.  As you will quickly
  notice, it is not really very close to the source.  This
  representation is not very useful to even sophisticated users.
  
\item The Virtual Machine (VM, or IR2) representation of the program.
  This dump represents the generated code as sequences of ``Virtual
  OPerations'' (VOPs.)  This representation is intermediate between
  the source and the assembly code\dash{}each VOP corresponds fairly
  directly to some primitive function or construct, but a given VOP
  also has a fairly predictable instruction sequence.  An operation
  (such as \code{+}) may have multiple implementations with different
  cost and applicability.  The choice of a particular VOP such as
  \code{+/fixnum} or \code{+/single-float} represents this choice of
  implementation.  Once you are familiar with it, the VM
  representation is probably the most useful for determining what
  implementation has been used.
  
\item An assembly listing, annotated with the VOP responsible for
  generating the instructions.  This listing is useful for figuring
  out what a VOP does and how it is implemented in a particular
  context, but its large size makes it more difficult to read.
  
\item A disassembly of the generated code, which has all
  pseudo-operations expanded out, but is not annotated with VOPs.
\end{itemize}


Note that trace file generation takes much space and time, since the trace file
is tens of times larger than the source file.  To avoid huge confusing trace
files and much wasted time, it is best to separate the critical program portion
into its own file and then generate the trace file from this small file.


\section{Efficiency Notes}
\label{efficiency-notes}
\cindex{efficiency notes}
\cpsubindex{notes}{efficiency}
\cindex{tuning}

Efficiency notes are messages that warn the user that the compiler has
chosen a relatively inefficient implementation for some operation.
Usually an efficiency note reflects the compiler's desire for more
type information.  If the type of the values concerned is known to the
programmer, then additional declarations can be used to get a more
efficient implementation.

Efficiency notes are controlled by the
\code{extensions:inhibit-warnings} (\pxlref{optimize-declaration})
optimization quality. When \code{speed} is greater than
\code{extensions:inhibit-warnings}, efficiency notes are enabled.
Note that this implicitly enables efficiency notes whenever
\code{speed} is increased from its default of \code{1}.

Consider this program with an obscure missing declaration:

\begin{lisp}
(defun eff-note (x y z)
  (declare (fixnum x y z))
  (the fixnum (+ x y z)))
\end{lisp}

If compiled with \code{\w{(speed 3) (safety 0)}}, this note is given:

\begin{example}
In: DEFUN EFF-NOTE
  (+ X Y Z)
==>
  (+ (+ X Y) Z)
Note: Forced to do inline (signed-byte 32) arithmetic (cost 3).
      Unable to do inline fixnum arithmetic (cost 2) because:
      The first argument is a (INTEGER -1073741824 1073741822),
      not a FIXNUM.
\end{example}

This efficiency note tells us that the result of the intermediate
computation \code{\w{(+ x y)}} is not known to be a \code{fixnum}, so
the addition of the intermediate sum to \code{z} must be done less
efficiently.  This can be fixed by changing the definition of
\code{eff-note}:

\begin{lisp}
(defun eff-note (x y z)
  (declare (fixnum x y z))
  (the fixnum (+ (the fixnum (+ x y)) z)))
\end{lisp}


\subsection{Type Uncertainty}
\cpsubindex{types}{uncertainty}
\cindex{uncertainty of types}

The main cause of inefficiency is the compiler's lack of adequate
information about the types of function argument and result values.
Many important operations (such as arithmetic) have an inefficient
general (generic) case, but have efficient implementations that can
usually be used if there is sufficient argument type information.

Type efficiency notes are given when a value's type is uncertain.
There is an important distinction between values that are {\em not
known} to be of a good type (uncertain) and values that are {\em known
not} to be of a good type. Efficiency notes are given mainly for the
first case (uncertain types.) If it is clear to the compiler that that
there is not an efficient implementation for a particular function
call, then an efficiency note will only be given if the
\code{extensions:inhibit-warnings} optimization quality is \code{0}
(\pxlref{optimize-declaration}.)

In other words, the default efficiency notes only suggest that you add
declarations, not that you change the semantics of your program so
that an efficient implementation will apply.  For example, compilation
of this form will not give an efficiency note:
\begin{lisp}
(elt (the list l) i)
\end{lisp}
even though a vector access is more efficient than indexing a list.


\subsection{Efficiency Notes and Type Checking}
\cpsubindex{type checking}{efficiency of}
\cpsubindex{efficiency}{of type checking}
\cpsubindex{optimization}{type check}

It is important that the \code{eff-note} example above used
\w{\code{(safety 0)}}.  When type checking is enabled, you may get apparently
spurious efficiency notes.  With \w{\code{(safety 1)}}, the note has this extra
line on the end:

\begin{example}
The result is a (INTEGER -1610612736 1610612733), not a FIXNUM.
\end{example}

This seems strange, since there is a \code{the} declaration on the result of that
second addition.

In fact, the inefficiency is real, and is a consequence of \python{}'s
treating declarations as assertions to be verified.  The compiler
can't assume that the result type declaration is true\dash{}it must
generate the result and then test whether it is of the appropriate
type.

In practice, this means that when you are tuning a program to run
without type checks, you should work from the efficiency notes
generated by unsafe compilation.  If you want code to run efficiently
with type checking, then you should pay attention to all the
efficiency notes that you get during safe compilation.  Since user
supplied output type assertions (e.g., from \code{the}) are
disregarded when selecting operation implementations for safe code,
you must somehow give the compiler information that allows it to prove
that the result truly must be of a good type.  In our example, it
could be done by constraining the argument types more:

\begin{lisp}
(defun eff-note (x y z)
  (declare (type (unsigned-byte 18) x y z))
  (+ x y z))
\end{lisp}

Of course, this declaration is acceptable only if the arguments to \code{eff-note}
always \var{are} \w{\code{(unsigned-byte 18)}} integers.


\subsection{Representation Efficiency Notes}
\label{representation-eff-note}
\cindex{representation efficiency notes}
\cpsubindex{efficiency notes}{for representation}
\cindex{object representation efficiency notes}
\cindex{stack numbers}
\cindex{non-descriptor representations}
\cpsubindex{descriptor representations}{forcing of}

When operating on values that have non-descriptor representations
(\pxlref{non-descriptor}), there can be a substantial time and consing
penalty for converting to and from descriptor representations.  For
this reason, the compiler gives an efficiency note whenever it is
forced to do a representation coercion more expensive than
\varref{efficiency-note-cost-threshold}.

Inefficient representation coercions may be due to type uncertainty,
as in this example:

\begin{lisp}
(defun set-flo (x)
  (declare (single-float x))
  (prog ((var 0.0))
    (setq var (gorp))
    (setq var x)
    (return var)))
\end{lisp}

which produces this efficiency note:

\begin{example}
In: DEFUN SET-FLO
  (SETQ VAR X)
Note: Doing float to pointer coercion (cost 13) from X to VAR.
\end{example}

The variable \code{var} is not known to always hold values of type
\code{single-float}, so a descriptor representation must be used for its value.
In this sort of situation, adding a declaration will eliminate the inefficiency.

Often inefficient representation conversions are not due to type
uncertainty\dash{}instead, they result from evaluating a
non-descriptor expression in a context that requires a descriptor
result:

\begin{itemize} 
\item Assignment to or initialization of any data structure other than
  a specialized array (\pxlref{specialized-array-types}), or
  
\item Assignment to a \code{special} variable, or
  
\item Passing as an argument or returning as a value in any function
  call that is not a local call (\pxlref{number-local-call}.)
\end{itemize}

If such inefficient coercions appear in a ``hot spot'' in the program, data
structures redesign or program reorganization may be necessary to improve
efficiency.  See sections \ref{block-compilation}, \ref{numeric-types} and
\ref{profiling}.

Because representation selection is done rather late in compilation,
the source context in these efficiency notes is somewhat vague, making
interpretation more difficult.  This is a fairly straightforward
example:

\begin{lisp}
(defun cf+ (x y)
  (declare (single-float x y))
  (cons (+ x y) t))
\end{lisp}

which gives this efficiency note:

\begin{example}
In: DEFUN CF+
  (CONS (+ X Y) T)
Note: Doing float to pointer coercion (cost 13), for:
      The first argument of CONS.
\end{example}

The source context form is almost always the form that receives the value being
coerced (as it is in the preceding example), but can also be the source form
which generates the coerced value.  Compiling this example:

\begin{lisp}
(defun if-cf+ (x y)
  (declare (single-float x y))
  (cons (if (grue) (+ x y) (snoc)) t))
\end{lisp}

produces this note:

\begin{example}
In: DEFUN IF-CF+
  (+ X Y)
Note: Doing float to pointer coercion (cost 13).
\end{example}

In either case, the note's text explanation attempts to include
additional information about what locations are the source and
destination of the coercion.  Here are some example notes:
\begin{example}
  (IF (GRUE) X (SNOC))
Note: Doing float to pointer coercion (cost 13) from X.

  (SETQ VAR X)
Note: Doing float to pointer coercion (cost 13) from X to VAR.
\end{example}
Note that the return value of a function is also a place to which coercions may
have to be done:
\begin{example}
  (DEFUN F+ (X Y) (DECLARE (SINGLE-FLOAT X Y)) (+ X Y))
Note: Doing float to pointer coercion (cost 13) to "<return value>".
\end{example}
Sometimes the compiler is unable to determine a name for the source or
destination, in which case the source context is the only clue.


\subsection{Verbosity Control}
\cpsubindex{verbosity}{of efficiency notes}
\cpsubindex{efficiency notes}{verbosity}

These variables control the verbosity of efficiency notes:

\begin{defvar}{}{efficiency-note-cost-threshold}
  
  Before printing some efficiency notes, the compiler compares the
  value of this variable to the difference in cost between the chosen
  implementation and the best potential implementation.  If the
  difference is not greater than this limit, then no note is printed.
  The units are implementation dependent; the initial value suppresses
  notes about ``trivial'' inefficiencies.  A value of \code{1} will
  note any inefficiency.
\end{defvar}

\begin{defvar}{}{efficiency-note-limit}
  
  When printing some efficiency notes, the compiler reports possible
  efficient implementations.  The initial value of \code{2} prevents
  excessively long efficiency notes in the common case where there is
  no type information, so all implementations are possible.
\end{defvar}


\section{Profiling}
\cindex{profiling}
\cindex{timing}
\cindex{consing}
\cindex{tuning}
\label{profiling}

The first step in improving a program's performance is to profile the
activity of the program to find where it spends its time.  The best
way to do this is to use the profiling utility found in the
\code{profile} package.  This package provides a macro \code{profile}
that encapsulates functions with statistics gathering code.


\subsection{Profile Interface}

\begin{defvar}{profile:}{timed-functions}
  
  This variable holds a list of all functions that are currently being
  profiled.
\end{defvar}

\begin{defmac}{profile:}{profile}{%
    \args{\mstar{\var{name} \mor \kwd{callers} \code{t}}}}
  
  This macro wraps profiling code around the named functions.  As in
  \code{trace}, the \var{name}s are not evaluated.  If a function is
  already profiled, then the function is unprofiled and reprofiled
  (useful to notice function redefinition.)  A warning is printed for
  each name that is not a defined function.
  
  If \kwd{callers \var{t}} is specified, then each function that calls
  this function is recorded along with the number of calls made.
\end{defmac}

\begin{defmac}{profile:}{unprofile}{%
    \args{\mstar{\var{name}}}}
  
  This macro removes profiling code from the named functions.  If no
  \var{name}s are supplied, all currently profiled functions are
  unprofiled.
\end{defmac}

\begin{defmac}{profile:}{profile-all}{%
    \args{\keys{\kwd{package} \kwd{callers-p}}}}
  
  This macro in effect calls \code{profile:profile} for each
  function in the specified package which defaults to
  \code{*package*}.  \kwd{callers-p} has the same meaning as in
  \code{profile:profile}.
\end{defmac}

\begin{defmac}{profile:}{report-time}{\args{\mstar{\var{name}}}}
  
  This macro prints a report for each \var{name}d function of the
  following information:
  \begin{itemize}
  \item The total CPU time used in that function for all calls,
  
  \item the total number of bytes consed in that function for all
    calls,
  
  \item the total number of calls,
  
  \item the average amount of CPU time per call.
  \end{itemize}
  Summary totals of the CPU time, consing and calls columns are
  printed.  An estimate of the profiling overhead is also printed (see
  below).  If no \var{name}s are supplied, then the times for all
  currently profiled functions are printed.
\end{defmac}

\begin{defmac}{}{reset-time}{\args{\mstar{\var{name}}}}
  
  This macro resets the profiling counters associated with the
  \var{name}d functions.  If no \var{name}s are supplied, then all
  currently profiled functions are reset.
\end{defmac}


\subsection{Profiling Techniques}

Start by profiling big pieces of a program, then carefully choose which
functions close to, but not in, the inner loop are to be profiled next.
Avoid profiling functions that are called by other profiled functions, since
this opens the possibility of profiling overhead being included in the reported
times.

If the per-call time reported is less than 1/10 second, then consider the clock
resolution and profiling overhead before you believe the time.  It may be that
you will need to run your program many times in order to average out to a
higher resolution.


\subsection{Nested or Recursive Calls}

The profiler attempts to compensate for nested or recursive calls.  Time and
consing overhead will be charged to the dynamically innermost (most recent)
call to a profiled function.  So profiling a subfunction of a profiled function
will cause the reported time for the outer function to decrease.  However if an
inner function has a large number of calls, some of the profiling overhead may
``leak'' into the reported time for the outer function.  In general, be wary of
profiling short functions that are called many times.


\subsection{Clock resolution}

Unless you are very lucky, the length of your machine's clock ``tick'' is
probably much longer than the time it takes simple function to run.  For
example, on the IBM RT, the clock resolution is 1/50 second.  This means that
if a function is only called a few times, then only the first couple decimal
places are really meaningful.  

Note however, that if a function is called many times, then the statistical
averaging across all calls should result in increased resolution.  For example,
on the IBM RT, if a function is called a thousand times, then a resolution of
tens of microseconds can be expected.

\subsection{Profiling overhead}

The added profiling code takes time to run every time that the profiled
function is called, which can disrupt the attempt to collect timing
information.  In order to avoid serious inflation of the times for functions
that take little time to run, an estimate of the overhead due to profiling is
subtracted from the times reported for each function.

Although this correction works fairly well, it is not totally accurate,
resulting in times that become increasingly meaningless for functions with
short runtimes.  This is only a concern when the estimated profiling overhead
is many times larger than reported total CPU time.

The estimated profiling overhead is not represented in the reported total CPU
time.  The sum of total CPU time and the estimated profiling overhead should be
close to the total CPU time for the entire profiling run (as determined by the
\code{time} macro.)  Time unaccounted for is probably being used by functions that
you forgot to profile.

\subsection{Additional Timing Utilities}

\begin{defmac}{}{time}{ \args{\var{form}}}

  This macro evaluates \var{form}, prints some timing and memory
  allocation information to \code{*trace-output*}, and returns any
  values that \var{form} returns.  The timing information includes
  real time, user run time, and system run time.  This macro executes
  a form and reports the time and consing overhead.  If the
  \code{time} form is not compiled (e.g. it was typed at top-level),
  then \code{compile} will be called on the form to give more accurate
  timing information.  If you really want to time interpreted speed,
  you can say:
\begin{lisp}
(time (eval '\var{form}))
\end{lisp}
Things that execute fairly quickly should be timed more than once,
since there may be more paging overhead in the first timing.  To
increase the accuracy of very short times, you can time multiple
evaluations:
\begin{lisp}
(time (dotimes (i 100) \var{form}))
\end{lisp}
\end{defmac}

\begin{defun}{extensions:}{get-bytes-consed}{}
  
  This function returns the number of bytes allocated since the first
  time you called it.  The first time it is called it returns zero.
  The above profiling routines use this to report consing information.
\end{defun}

\begin{defvar}{extensions:}{gc-run-time}
  
  This variable accumulates the run-time consumed by garbage
  collection, in the units returned by
  \findexed{get-internal-run-time}.
\end{defvar}

\begin{defconst}{}{internal-time-units-per-second}
The value of internal-time-units-per-second is 100.
\end{defconst}

\subsection{A Note on Timing}
\cpsubindex{CPU time}{interpretation of}
\cpsubindex{run time}{interpretation of}
\cindex{interpretation of run time}

There are two general kinds of timing information provided by the
\code{time} macro and other profiling utilities: real time and run
time.  Real time is elapsed, wall clock time.  It will be affected in
a fairly obvious way by any other activity on the machine.  The more
other processes contending for CPU and memory, the more real time will
increase.  This means that real time measurements are difficult to
replicate, though this is less true on a dedicated workstation.  The
advantage of real time is that it is real.  It tells you really how
long the program took to run under the benchmarking conditions.  The
problem is that you don't know exactly what those conditions were.

Run time is the amount of time that the processor supposedly spent
running the program, as opposed to waiting for I/O or running other
processes.  ``User run time'' and ``system run time'' are numbers
reported by the Unix kernel.  They are supposed to be a measure of how
much time the processor spent running your ``user'' program (which
will include GC overhead, etc.), and the amount of time that the
kernel spent running ``on your behalf.''

Ideally, user time should be totally unaffected by benchmarking
conditions; in reality user time does depend on other system activity,
though in rather non-obvious ways.

System time will clearly depend on benchmarking conditions.  In Lisp
benchmarking, paging activity increases system run time (but not by as much
as it increases real time, since the kernel spends some time waiting for
the disk, and this is not run time, kernel or otherwise.)

In my experience, the biggest trap in interpreting kernel/user run time is
to look only at user time.  In reality, it seems that the \var{sum} of kernel
and user time is more reproducible.  The problem is that as system activity
increases, there is a spurious \var{decrease} in user run time.  In effect, as
paging, etc., increases, user time leaks into system time.

So, in practice, the only way to get truly reproducible results is to run
with the same competing activity on the system.  Try to run on a machine
with nobody else logged in, and check with ``ps aux'' to see if there are any
system processes munching large amounts of CPU or memory.  If the ratio
between real time and the sum of user and system time varies much between
runs, then you have a problem.


\subsection{Benchmarking Techniques}
\cindex{benchmarking techniques}

Given these imperfect timing tools, how do should you do benchmarking?  The
answer depends on whether you are trying to measure improvements in the
performance of a single program on the same hardware, or if you are trying to
compare the performance of different programs and/or different hardware.

For the first use (measuring the effect of program modifications with
constant hardware), you should look at \var{both} system+user and real time to
understand what effect the change had on CPU use, and on I/O (including
paging.)  If you are working on a CPU intensive program, the change in
system+user time will give you a moderately reproducible measure of
performance across a fairly wide range of system conditions.  For a CPU
intensive program, you can think of system+user as ``how long it would have
taken to run if I had my own machine.''  So in the case of comparing CPU
intensive programs, system+user time is relatively real, and reasonable to
use.

For programs that spend a substantial amount of their time paging, you
really can't predict elapsed time under a given operating condition without
benchmarking in that condition.  User or system+user time may be fairly
reproducible, but it is also relatively meaningless, since in a paging or
I/O intensive program, the program is spending its time waiting, not
running, and system time and user time are both measures of run time.
A change that reduces run time might increase real time by increasing
paging.

Another common use for benchmarking is comparing the performance of
the same program on different hardware.  You want to know which
machine to run your program on.  For comparing different machines
(operating systems, etc.), the only way to compare that makes sense is
to set up the machines in \var{exactly} the way that they will
\var{normally} be run, and then measure \var{real} time.  If the
program will normally be run along with X, then run X.  If the program
will normally be run on a dedicated workstation, then be sure nobody
else is on the benchmarking machine.  If the program will normally be
run on a machine with three other Lisp jobs, then run three other Lisp
jobs.  If the program will normally be run on a machine with 64MB of
memory, then run with 64MB.  Here, ``normal'' means ``normal for that
machine''.  

If you have a program you believe to be CPU intensive, then you might be
tempted to compare ``run'' times across systems, hoping to get a meaningful
result even if the benchmarking isn't done under the expected running
condition.  Don't to this, for two reasons:

\begin{itemize}  
\item The operating systems might not compute run time in the same
  way.
  
\item Under the real running condition, the program might not be CPU
  intensive after all.
\end{itemize}


In the end, only real time means anything\dash{}it is the amount of time you
have to wait for the result.  The only valid uses for run time are:

\begin{itemize}
\item To develop insight into the program.  For example, if run time
  is much less than elapsed time, then you are probably spending lots
  of time paging.
  
\item To evaluate the relative performance of CPU intensive programs
  in the same environment.
\end{itemize}
